\begin{afterword}
\summaryhead

The post begins from beliefs whose flaws a ten-year-old could see if hearing them for the first time. Even so, a Nobel laureate who knows Bayesian mathematics, Robert Aumann, is an Orthodox Jew. The author feels confident that Aumann must have questioned his faith, and says he did not doubt successfully. So anyone may hold a belief as bad as religion without knowing it.

The hard question is not how to reject a false belief but how to tell whether a belief is false. Reciting standard replies, including the rationalist's ``Occam's Razor!'', is the habit that keeps believers in their religion. A true crisis of faith, one that could go either way, needs the desperate effort a theist would need to leave a religion: block the cached replies, imagine what a skeptic would say to your answer, and think the thought that hurts most.

Not every belief needs this. The warning signs are a belief held long, surrounded by known arguments, costly to give up, emotionally charged and part of one's personality; they apply to Dawkins's evolutionary biology as much as to the Pope's Catholicism. For such a belief, rest the day before, set aside quiet, uninterrupted hours, clear your mind of standard arguments and make a desperate effort to doubt. The post lists earlier essays that supply the technique's parts, and ends with isshokenmei, the lifelong effort to be rational, of which the crisis is the critical point.

\responsehead

The post's central point is sound. It applies its test to its readers' own side: the rationalist who answers ``Occam's Razor!'' is reciting a standard reply, as the theist is, and its warning signs are stated for Dawkins as well as the Pope. Its reframing of the problem, from how to reject a false belief to how to tell whether a belief is false, is the right question. The problems are the evidence it offers that the danger is real, and the technique it offers to meet it.

The evidence is one man, and the post does not cite what he had said. That Aumann questioned his faith is offered as a guess. In a 2005 interview Aumann described his religion as ``mainly an emotional and aesthetic'' experience, ``not about whether the earth is 5,765 years old,'' and nearly independent of science, while saying belief is an important part of it. Where a religious claim could be tested (the Bible codes, whose study he had promoted), a committee he chaired failed to replicate a result, and he called the codes ``improbable.'' More basically, his keeping his faith counts as a failed doubt only if the faith is false. The post asserts that it is, in its first paragraph. So the case shows the danger only to readers who already share the verdict, while the post's own question is how one would know. The post it links on doubt counts a doubt that ends with the belief kept as ``a resolution.''

The technique is advice without a test. The post describes no crisis, the author's or anyone's, and gives no evidence that rest, quiet hours and a cleared mind change the outcome. Its setting is close to Descartes's in the First Meditation, which the post does not mention.

Its footnote applies the rule only to the author's opponents. The people it tells to question their formulation of Occam's razor are those who think it rules out many-worlds; the author's own case for many-worlds was based on a formulation of Occam's razor, and the footnote does not question it.

Two smaller points. ``How dare you believe anything'' is narrowed three paragraphs later to beliefs with warning signs. And the ``camel has two humps'' that the author wonders about probably alludes to a paper whose co-author later retracted its strongest claims.

In short: a post that rightly asks rationalists to doubt their own standard replies, but offers as evidence of the danger one scientist whose own account it does not cite and whose failure it assumes, and offers a technique it does not show at work.
\end{afterword}
